\section{Introduction}\label{sec:introduction}

The word problem for a group with a specified finite generating set asks
whether a given word in the generators and their inverses represents the
identity. It is \emph{decidable} if an algorithm answers this question for
every word. A group is \emph{simple} here if it is nontrivial and its only
normal subgroups are the identity subgroup and the whole group.

\begin{theorem}\label{thm:main}
For every finitely generated group \(G\), the following are equivalent:
\begin{enumerate}[label=\textup{(\roman*)}]
\item \(G\) has decidable word problem.
\item \(G\) embeds in a finitely presented simple group.
\end{enumerate}
\end{theorem}

An embedding throughout means an injective group homomorphism. In the
constructive implication, \(G\) is supplied with a finite generating set
and a word-problem decider; it need not be finitely presented.
Theorem~\ref{thm:main} resolves the Boone--Higman conjecture positively.
